\SourceLeaf{173}
\Article{36}
Another part of our theory open to doubt is the assumption that the agents of the Earth's magnetic force are situated exclusively in its interior.

If the immediate causes are to be sought wholly or partly outside the Earth, we can. if we rule out groundless speculation and confine ourselves to what is scientifically known, think only of galvanic currents. The air of the atmosphere does not conduct such currents, and neither does empty space. Our knowledge therefore fails us when we look for a medium to carry galvanic currents in the upper regions. Yet the enigmatic phenomena of the northern lights, in which electricity in motion appears to play a leading role, prevent us from denying outright the possibility of such currents on the strength of that ignorance alone. It remains interesting in any case to investigate what magnetic effect they would produce on the Earth's surface.

\Article{37}
Let us therefore suppose that steady galvanic currents occur in a region $S$ that encloses the Earth like a vault or a shell. Let us denote the whole region enclosed by $S$ by $S'$, and the whole exterior region that encloses $S$ and $S'$ by $S''$. Whatever the configuration of those galvanic currents, they can always be replaced by a fictitious distribution of magnetic fluids within the region $S$ itself, which exerts exactly the same magnetic effect throughout the rest of space, $S'$ and $S''$, as the currents do. This important proposition, already mentioned in Article 3, rests first on the fact that those currents can be decomposed into infinitely many elementary currents (that is, currents that may be treated as linear). It rests second on the well-known theorem, first proved, to my knowledge, by \textsc{Ampère}, that any linear current bounding an arbitrary surface can be replaced, with the stated effect, by a distribution of the magnetic fluids on the two sides of that surface at immeasurably small distances from it. It rests third on the evident possibility of assigning to every closed line lying within $S$ a surface bounded by it that also lies entirely within $S$.

\SourceLeaf{174}
Now let $-v$ denote the sum of all the quotients obtained by dividing each element of that fictitious magnetic fluid by its distance from an unspecified point $O$ in $S'$ or $S''$. Here the elements of the southern fluid must, of course, be taken as negative. The partial derivatives of $v$ (just like those of $V$ in our theory above) then express the components of the magnetic force produced at $O$ by the galvanic currents.

\Article{38}
Although a detailed development of the theory from which the proposition used in the preceding article is taken must be left for another occasion, one important point concerning it still deserves mention here. Suppose that two \textit{different} surfaces $F$, $F'$ are constructed, each bounded by the same linear current $G$. For brevity, consider only the simplest case, in which they have no point in common apart from their common boundary line. The surfaces then enclose a three-dimensional region. If $O$ lies outside this region, we obtain \textit{one and the same} value for the component of $v$ associated with $G$, whether we distribute the magnetic fluids on $F$ or on $F'$. This value equals the product of the intensity of the galvanic current $G$ (measured in a suitable unit) and the solid angle with its vertex at $O$ and bounded by the straight lines drawn from $O$ to the points of $G$. Equivalently, it equals that intensity multiplied by the portion of the spherical surface of radius 1 about $O$ that is the common projection of both $F$ and $F'$. If, however, $O$ lies inside the region enclosed by $F$ and $F'$, the two values of the part of $v$ in question differ according to whether the magnetic fluids are distributed on $F$ or on $F'$. This is because the values correspond to different portions of the spherical surface just mentioned, which together make up the entire surface. But since the direction of the galvanic current is opposite relative to $F$ and $F'$, opposite signs must be assigned to the current's intensity in the two cases when multiplying by the portions of the spherical surface. It follows that the algebraic difference between the two values
\SourceLeaf{175}
of the part of $v$ in question equals the current's intensity multiplied by the entire spherical surface, or by $4\pi$.

We readily conclude that if $O$ lies in $S''$, the value of $v$ is entirely independent of the choice of connecting surfaces. If, however, $O$ lies in $S'$, the absolute value of $v$ does depend on that choice, while the differentials of $v$ are independent of it.

The very fruitful theorem touched on here still requires further explanation in its full generality. According to this theorem, for the magnetic effect of a linear galvanic current, the product of its intensity and the portion of the spherical surface bounded by the projection of the current line from $O$ has the same significance as the elements of mass divided by their distance from $O$ have for attractive or repulsive forces. Those explanations must be reserved for a detailed treatment of the subject.
